\documentclass[11pt]{article}
\usepackage[T1]{fontenc}
\usepackage[utf8]{inputenc}
\usepackage{lmodern}
\usepackage[a4paper,margin=1in]{geometry}
\usepackage{microtype}
\usepackage{amsmath,amssymb,amsthm,mathtools}
\usepackage{booktabs,array}
\usepackage[dvipsnames]{xcolor}
\usepackage[colorlinks=true,linkcolor=MidnightBlue,citecolor=MidnightBlue,urlcolor=MidnightBlue]{hyperref}
\usepackage{enumitem}
\usepackage{setspace}
\usepackage{fancyhdr}
\usepackage{titlesec}
\setlength{\headheight}{14pt}
\pagestyle{fancy}
\fancyhf{}
\fancyhead[L]{Round-7 audit / response (Claude)}
\fancyhead[R]{May 2026}
\fancyfoot[C]{\thepage}
\titleformat{\section}{\large\bfseries\color{MidnightBlue}}{\thesection}{0.6em}{}
\titleformat{\subsection}{\normalsize\bfseries\color{MidnightBlue}}{\thesubsection}{0.5em}{}
\setlist[itemize]{topsep=4pt,itemsep=2pt,leftmargin=1.5em}
\setlist[enumerate]{topsep=4pt,itemsep=2pt,leftmargin=1.8em}
\onehalfspacing

\newtheorem{theorem}{Theorem}[section]
\newtheorem{proposition}[theorem]{Proposition}
\newtheorem{lemma}[theorem]{Lemma}
\newtheorem{corollary}[theorem]{Corollary}
\theoremstyle{definition}
\newtheorem{definition}[theorem]{Definition}
\theoremstyle{remark}
\newtheorem{remark}[theorem]{Remark}
\newtheorem{example}[theorem]{Example}

\newcommand{\DR}{\mathsf{DR}}
\newcommand{\PO}{\mathsf{PO}}
\newcommand{\PP}{\mathsf{PP}}
\newcommand{\PPI}{\mathsf{PPI}}
\newcommand{\EQ}{\mathsf{EQ}}
\newcommand{\Base}{\mathsf{B}}
\newcommand{\Inv}{\operatorname{Inv}}
\newcommand{\comp}{\operatorname{comp}}
\newcommand{\ALCIRCC}[1]{\mathcal{ALCI}_{\mathsf{RCC#1}}}
\newcommand{\Q}{\mathcal Q}
\newcommand{\U}{\mathcal U}
\newcommand{\T}{\mathcal T}
\newcommand{\Pair}{\mathsf{Pair}}
\newcommand{\Tri}{\mathsf{Tri}}
\newcommand{\Pos}{\mathsf{Pos}}
\newcommand{\Inc}{\mathsf{Inc}}
\newcommand{\Port}{\mathsf{Port}}
\newcommand{\Ctx}{\mathsf{Ctx}}
\newcommand{\Mos}{\mathsf{Mos}}
\newcommand{\Side}{\mathsf{Side}}
\newcommand{\Vrt}{\mathsf{Vert}}
\newcommand{\lab}{\operatorname{lab}}
\newcommand{\Stratum}{\mathsf{Stratum}}
\newcommand{\Subsrc}{\mathsf{Subsrc}}
\newcommand{\eqc}{\equiv_{\EQ}}
\newcommand{\Cl}{\operatorname{cl}}
\newcommand{\self}{\mathsf{self}}
\newcommand{\eqtag}{\mathsf{eq}}
\newcommand{\up}{\mathsf{up}}
\newcommand{\down}{\mathsf{down}}
\newcommand{\side}{\mathsf{side}}
\newcommand{\front}{\mathsf{front}}

\begin{document}

\begin{center}
{\LARGE\bfseries Round-7 Audit and Response (Claude)\par}
\vspace{0.5em}
{\large Adopting GPT-5.5's pair-shape / incidence-tag repair as the
canonical statement\par}
\vspace{0.75em}
{\normalsize Claude (Anthropic, Opus~4.7), prompted by Michael Wessel}\par
\vspace{0.45em}
{\small May~2026}
\end{center}

\begin{center}
\fbox{\parbox{0.92\textwidth}{%
\textbf{Disclaimer.}\;
This document was drafted by Claude (Anthropic), prompted by Michael
Wessel.  The results have not been independently verified by human
domain experts.  We invite scrutiny and corrections.}}
\end{center}

\begin{abstract}
GPT-5.5's round-7 review of Claude's round-6 internal-coherence pass
of the consolidated v2 proof identifies one central architectural
defect in Claude's manuscript: the abstract pair-label function
$L_\Q:\Pos_\Q^2\to\Base$ with diagonal $L_\Q(\pi,\pi)=\EQ$ collapses
distinct laps of a blocking cycle into semantic equality, because in
a finite regular certificate two fresh occurrences $u_i,u_j$ of a
one-state upward cycle share the same abstract position symbol $\pi$
and would be assigned $\lab(u_i,u_j)=L_\Q(\pi,\pi)=\EQ$.  This breaks
the canonical certificate for $\exists\PP.\top\sqcap\forall\PP.\exists\PP.\top$
and undermines the central blocking-not-equality slogan.  We accept
that the round-6 architecture cannot be locally patched: the
diagonal-$\EQ$ collapse is intrinsic to indexing labels by abstract
position symbol pairs.  We therefore adopt GPT-5.5's round-7 proof
sketch~\cite{GPT55Round7Sketch} as the new canonical statement; the
expanded full-details version~\cite{GPT55Round7Expanded} is its
detailed companion; the alternative parity-tree-automata
proof~\cite{GPT55AutomataRoute} is retained as an independent
witness.  Claude's role shifts back to verification and audit.

This response paper (i) summarizes the round-7 architectural move
(incidence tags and pair shapes), (ii) maps the six gaps (G1)--(G6)
of Claude's round-6 paper and the five gaps (D1)--(D5) of GPT-5.5's
round-7 review to their resolution in the round-7 sketch, (iii)
reports three self-contained Python verifications
(\texttt{wp7\_selfcontained\_side\_witness.py},
\texttt{wp8\_round7\_blocking\_chain.py},
\texttt{wp9\_round7\_split\_copies.py}) corroborating the three
central stress concepts of the round-7 architecture, and (iv)
discusses the no-automata route as the canonical path with the
automata route as the independent witness.
\end{abstract}

\tableofcontents
\bigskip

%% ==================================================================
\section{Why round-6 cannot be locally patched}
\label{sec:round6-cannot-be-patched}
%% ==================================================================

GPT-5.5's round-7 review~\cite{GPT55Round7Review} identifies a single
critical defect in Claude's round-6 manuscript that is independent of
the various surface inconsistencies.  Round-6 defined the concrete
label on every unfolding pair by
\[
  \lab(u,v)\;:=\;L_\Q(q(u),q(v)),
\]
where $q:\Omega(\U(\Q))\to\Pos_\Q$ sends each concrete occurrence to
an abstract position symbol and $L_\Q:\Pos_\Q^2\to\Base$ is the
declared abstract pair-label function with the standard reflexivity
clause $L_\Q(\pi,\pi)=\EQ$.

\paragraph{The collapse.}
In a finite regular certificate, the abstract position alphabet is
finite while the unfolding is infinite.  Blocking deliberately reuses
the same abstract position symbol on infinitely many fresh semantic
occurrences.  Take the canonical one-state upward cycle for
$C_\uparrow:=\exists\PP.\top\sqcap\forall\PP.\exists\PP.\top$, whose
unfolding is
\[
  u_0\PP u_1\PP u_2\PP\cdots
\]
with all $u_i$ at the same finite profile, port, context, and stratum.
Then $q(u_i)=q(u_j)=\pi$ for every $i,j$, and the round-6 rule gives
\[
  \lab(u_i,u_j)\;=\;L_\Q(\pi,\pi)\;=\;\EQ,
\]
\emph{even for $i\ne j$}.  Two contradictory conclusions follow.
Either the quotient identifies $u_i$ and $u_j$, collapsing the strict
$\PP$ chain into a $\PP$-self-loop (which violates strong-$\EQ$
semantics on $\PP$); or the quotient does not identify them, in which
case two distinct elements of the quotient frame are $\EQ$-related
(which also violates strong-$\EQ$ semantics).  Either way the round-6
labelling does not represent the regular certificate it was designed
to represent.

\paragraph{Why no local patch suffices.}
A weaker version of the diagonal clause does not work.  Without
$L_\Q(\pi,\pi)=\EQ$, the reflexive case $u=v$ has no defined label.
With a side clause ``$L_\Q(\pi,\pi)=\EQ$ \emph{except when the two
occurrences are different}'', the function is no longer a function of
$\Pos_\Q^2$ alone:\@ to evaluate at $(q(u),q(v))=(\pi,\pi)$ one must
also know whether $u=v$, which is exactly the information $L_\Q$ over
position-symbol pairs cannot carry.  The defect is intrinsic to the
indexing.

\paragraph{The structural move.}
GPT-5.5's round-7 fix is to redefine the label function on a different
finite domain: \emph{occurrence-sensitive pair shapes}.  A pair shape
$\alpha(u,v)$ carries an explicit incidence tag $\iota(u,v)$ that
distinguishes literal identity ($u=v$, tag $\self$) from same-symbol
fresh copies (tag $\up$, $\down$, $\side_R$, or $\front_R$).  The
label function is
\[
  \ell_\Q:\Pair_\Q\to\Base,\qquad
  \ell_\Q(\self)=\EQ,\quad \ell_\Q(\eqtag)=\EQ,\quad
  \ell_\Q(\up)=\PP,\quad \ell_\Q(\down)=\PPI,\quad
  \ell_\Q(\side_R)=\ell_\Q(\front_R)=R.
\]
Reflexivity gives $\EQ$ only when the incidence tag is $\self$; two
fresh laps with $\iota(u_i,u_j)=\up$ get label $\PP$, never $\EQ$.

We accept this fix.  It is not a small revision of the round-6
formulation, because the entire round-6 architecture (the abstract
position symbol $\Pos_\Q$, the (V11) composition table on $\T_\Q$, the
(SC4$'$) per-position support closure, the round-5 (M6$'$) stratum
forest, and the round-6 (V15) cross-pair universal safety) is built
around the per-position labelling.  Replacing it requires reworking
the validity-clause numbering, the patchwork proof, the soundness
chain, and the verification artefacts.  We adopt the round-7 sketch
as the canonical proof manuscript and rebuild the verification layer
around the new architecture.

%% ==================================================================
\section{The round-7 architecture: incidence tags and pair shapes}
\label{sec:round7-architecture}
%% ==================================================================

We summarize the round-7 architecture from~\cite{GPT55Round7Sketch}
for self-containment.

\subsection{Incidence tags}

\begin{definition}[Incidence tags, after \cite{GPT55Round7Sketch} Def.~5.3]
\label{def:incidence}
For two concrete occurrences $u,v$ in the unfolding $\U(\Q)$, the
\emph{incidence tag} $\iota(u,v)\in\Inc$ takes one of the following
finite values:
\[
\begin{array}{ll}
\self  & u=v \text{ as concrete occurrences}\\
\eqtag & u\ne v \text{ and } u\eqc v \text{ by explicit equality ports}\\
\up    & v \text{ is a strict equality-aware proper superpart of } u\\
\down  & v \text{ is a strict equality-aware proper part of } u\\
\side_R & \text{the pair is named in a side mosaic with label }R\in\Base\\
\front_R & \text{the pair is residual sibling-frontier with }R\in\{\DR,\PO\}.
\end{array}
\]
The tags $\self$ and $\eqtag$ are \emph{occurrence-sensitive}: repeated
laps of a blocking cycle never receive $\self$ merely because their
finite profiles agree, and never receive $\eqtag$ unless explicitly
linked by the equality-port relation.
\end{definition}

\subsection{Pair shapes and the label function}

\begin{definition}[Pair shape and label, after \cite{GPT55Round7Sketch} Def.~5.4--5.5]
A \emph{pair shape} is a tuple
\[
  \alpha(u,v)\;=\;\bigl(s_u,p_u,s_v,p_v,\iota(u,v)\bigr),
\]
with $s_u,s_v$ profiles, $p_u\in\Port(s_u)$, $p_v\in\Port(s_v)$, and
$\iota(u,v)\in\Inc$.  The finite set of permitted pair shapes is
$\Pair_\Q$.  The label function is
\[
  \ell_\Q:\Pair_\Q\to\Base,
\]
with mandatory values
\[
  \ell_\Q(\self)=\EQ,\quad
  \ell_\Q(\eqtag)=\EQ,\quad
  \ell_\Q(\up)=\PP,\quad
  \ell_\Q(\down)=\PPI,\quad
  \ell_\Q(\side_R)=\ell_\Q(\front_R)=R.
\]
(For the $\self$/$\eqtag$ clauses, the profile/port fields are
determined by $u$ and $v$; for the others, the tag determines the
relation directly.)
\end{definition}

\begin{remark}\label{rem:no-collapse}
The diagonal collapse of round-6 is averted by Definition~\ref{def:incidence}.
Two distinct laps $u_i,u_j$ ($i\ne j$) of a blocking cycle satisfy
$u_i\ne u_j$ and (typically) $u_i\not\eqc u_j$.  Therefore
$\iota(u_i,u_j)\in\{\up,\down,\side_R,\front_R\}$, never $\self$ or
$\eqtag$.  In particular, $\ell_\Q(\alpha(u_i,u_j))\ne\EQ$ for any
distinct laps; for the one-state upward cycle, the label is $\PP$.
\end{remark}

\subsection{Triple shapes and composition (V6)}

\begin{definition}[Triple shape, after \cite{GPT55Round7Sketch} Def.~5.6]
A \emph{triple shape} is a tuple $\Theta=(\alpha_{12},\alpha_{23},\alpha_{13})$
of pair shapes satisfying coherence constraints: equality endpoints
must collapse the third pair label across $\eqc$; vertical chains must
project to vertical labels at the third edge; etc.  The finite set of
coherent triple shapes is $\Tri_\Q$.
\end{definition}

\paragraph{(V6) Triple composition.}
For every $\Theta=(\alpha_{12},\alpha_{23},\alpha_{13})\in\Tri_\Q$,
\[
  \ell_\Q(\alpha_{13})\;\in\;\comp\bigl(\ell_\Q(\alpha_{12}),\ell_\Q(\alpha_{23})\bigr).
\]
This is the round-7 successor of Claude's round-6 (V11).  The
difference is that triples are over \emph{pair shapes}, not over
position-symbol pairs, so the diagonal-$\EQ$ collapse cannot occur in
the composition check.

\subsection{Validity clauses (V1)--(V12)}

The round-7 validity clauses, summarized:

\begin{enumerate}[label=(V\arabic*)]
\item Type legality (Hintikka types).
\item Context legality (mosaics satisfy local axioms; local
$\PP/\PPI$ side labels are verticalized).
\item Pair-shape totality (every pair has exactly one shape).
\item Occurrence-sensitive equality ($\self$ only for $u=v$;
$\eqtag$ only for explicit equality-port pairs; no strict vertical
pair is $\eqtag$).
\item Typed equality congruence ($\eqtag$-mates have identical types
and identical $\ell_\Q$-labels to every shape).
\item Triple composition over $\Tri_\Q$ (above).
\item Vertical existential support ($\PP/\PPI$): equality-aware
strict reachability to a witness profile, with strict-later occurrences
on cyclic witnesses.
\item Vertical universal safety ($\PP/\PPI$): every strictly
reachable target has the witnessed type.
\item Side existential support ($\DR/\PO$): named mosaic position.
\item Side universal safety ($\DR/\PO$): every named side/frontier
target carries the type.
\item Request-closed cycles: every cycle is half-open
request-closed; cyclic witnesses are strict later occurrences in the
same or next lap.
\item Cycle universal safety: universal safety holds across the
entire cyclic order including wrap-around.
\end{enumerate}

Twelve clauses versus Claude's round-5/round-6 fifteen.  The
reduction is real: round-5 (V11) and round-6 (V15) are absorbed by
the round-7 (V6); round-5 (V12)/(V13)/(V14) are absorbed by the new
(V9)/(V5)/(V11) respectively (joint equality is folded into typed
congruence, declared overlap labels into mosaic legality, uniform
cycle representatives into the request-closed-cycle clause directly).

%% ==================================================================
\section{Mapping round-6 (G1)--(G6) gaps to round-7 resolutions}
\label{sec:round6-gaps}
%% ==================================================================

GPT-5.5's round-6 review listed six internal-coherence gaps in
Claude's round-6 consolidated paper.  All six are resolved by the
round-7 architectural move.

\begin{description}[leftmargin=2em]
\item[\textbf{G1 (flat side frontier).}]  Round-6 inlined (M6$'$) into
the Mosaic definition but left the round-4 ``Side-interface legality''
lemma asserting flat-frontier $\{\DR,\PO,\EQ\}$ on sibling positions.
Round-7 ((V2), Mosaic axiom (M6)) requires that any $\PP/\PPI$ side
label live inside a declared local vertical stratum; residual frontier
pairs are then restricted to $\{\DR,\PO\}$.  The flat-frontier lemma
is dropped.

\item[\textbf{G2 (SC4$'$ per-position vs.\ per-triple).}]  Round-6
weakened (SC4) to per-position but kept Quotient lifting and
Patchwork invoking per-triple mosaic coverage.  Round-7's
per-position support closure is consistent because composition is
checked globally on the triple-shape set $\Tri_\Q$ via (V6), and
patching is by the occurrence-sensitive triple-shape table
(\cite{GPT55Round7Sketch} Lemma 4.3), not by a per-triple mosaic.

\item[\textbf{G3 (cross-pair universal safety).}]  Round-6 added (V15)
to fix the under-proved universal-propagation lemma.  Round-7 folds
this into (V5) (typed equality congruence on mosaic positions) and
(V8)/(V10) (universal safety on vertical and side targets).  Cross-pair
universal safety is now a corollary of the universal-safety clauses on
the cycle and on the support mosaics.

\item[\textbf{G4 (extended-profile lasso, half-open cycle).}]  Round-7
adopts the half-open $[i,j)$ convention throughout and uses the
extended profile (including pending requests) for the request-closed
cycle lemma.  No tension between $\Pi$ and $\widehat\Pi$ remains.

\item[\textbf{G5 (bound consistency, V1--V15, stratum-forest grammar).}]
Round-7 has only 12 validity clauses and does not assert a polynomial
side-width.  The grammar enumerates over occurrence-sensitive pair
shapes and triple shapes; the bound on mosaic size is a crude
polynomial in $|\Cl(C_0)|$ ($p(n)=1+4n+3(4n+1)^3$ in
\cite{GPT55Round7Sketch}), sufficient for decidability.  Inconsistency
between cubic and recurrence bounds is moot.

\item[\textbf{G6 (self-contained verification script).}]  Round-7
verification: see \S\ref{sec:verification} below.  Three pure-stdlib
scripts (\texttt{wp7\_selfcontained\_side\_witness.py},
\texttt{wp8\_round7\_blocking\_chain.py},
\texttt{wp9\_round7\_split\_copies.py}) replace the
project-imports-dependent original.
\end{description}

%% ==================================================================
\section{Mapping round-7 (D1)--(D5) defects to their resolutions}
\label{sec:round7-defects}
%% ==================================================================

GPT-5.5's round-7 review of Claude's round-6 paper lists, in addition
to the central defect (\S\ref{sec:round6-cannot-be-patched}), five
further issues.  All are addressed by the round-7 architectural move.

\begin{description}[leftmargin=2em]
\item[\textbf{D1 (equality-port lift ambiguous).}]  Round-6 wrote
$u\eqc v$ iff $\Port(u)\eqc^\Q\Port(v)$ in the same generated
context or transitively across explicit links.  Read literally, this
identifies every occurrence using the same finite port.  Round-7
treats explicit equality-port links as \emph{occurrence-template} data
generating instantiated equality edges inside context instances; the
lift is $u\eqc v$ iff $u=v$ or there is a path of instantiated
explicit equality edges between $u$ and $v$.  Reflexivity at the
finite port does not collapse fresh blocked copies.

\item[\textbf{D2 (V9.b at the finite state-port level).}]  Round-7
restates V9.b at the occurrence-template level: it forbids explicit
equality between two positions that can be strict-vertical comparable
\emph{in the same unfolding instance}, without rejecting self-loop
blocking profiles by reflexivity.

\item[\textbf{D3 (stratum identity in abstract positions).}]  Round-6
kept a binary $\zeta\in\{\mathrm{vert},\mathrm{sib}\}$ field on
abstract positions, insufficient for nested strata.  Round-7's mosaic
declares vertical strata as a finite family $\mathcal V$, and the
co-vertical pair shape carries the stratum identity directly; the
binary flag is dropped.

\item[\textbf{D4 (recursive side-interface underspecified for chains).}]
Round-6's recursive side-interface construction was ambiguous for
side-witness chains $w_0\PP w_1\PP w_2$.  Round-7 (M6) is the simpler
condition: every $\PP/\PPI$ side label in a mosaic lies in a declared
vertical stratum; the stratum carries an explicit local containment
order, so chains are represented as ordered sequences inside a single
declared stratum.

\item[\textbf{D5 (universal wrap-around in pumped cycles).}]  Round-6's
request-cycle proof preserved existentials but the wrap-around case
of universals was under-proved.  Round-7 adds (V12) Cycle universal
safety: every $\forall R.D$ at any cycle position holds at every later
cycle position including wrap-around.  This is a finite cycle-word
check.
\end{description}

%% ==================================================================
\section{Verification corroboration}
\label{sec:verification}
%% ==================================================================

We supply three self-contained Python verifications, each in
\texttt{verification/python/}, that exercise the three central stress
concepts of the round-7 architecture.  Each script uses only the
Python standard library and runs from the verification archive alone.

\subsection{WP7: side-witness comparability}
\label{ssec:wp7}

\texttt{wp7\_selfcontained\_side\_witness.py} (round-6 (G6) repair)
builds the three-element abstract RCC5 frame $\Delta=\{x,y,z\}$ with
$\rho(x,y)=\DR$, $\rho(x,z)=\DR$, $\rho(y,z)=\PP$, $A=\{y\}$, $B=\{z\}$,
and verifies:
\begin{itemize}
\item JEPD, inverse, strong-$\EQ$, and the full RCC5 composition table
on every pair and triple;
\item $C_{\rm side}:=\exists\DR.(A\sqcap\exists\PP.B)\sqcap\exists\DR.B$
is SAT at $x$.
\end{itemize}
The script confirms the round-5/round-7 claim that two $\DR$-witnesses
of a source can themselves be $\PP/\PPI$-comparable.  The round-7
architecture handles this via local vertical stratification inside a
side mosaic.

\subsection{WP8: blocking is not equality}
\label{ssec:wp8}

\texttt{wp8\_round7\_blocking\_chain.py} (round-7 central repair)
truncates the canonical one-state upward cycle of
$C_\uparrow:=\exists\PP.\top\sqcap\forall\PP.\exists\PP.\top$ to a
five-element frame $\{u_0,\ldots,u_4\}$ with
$\rho(u_i,u_j)=\PP$ for $i<j$, $\PPI$ for $i>j$, $\EQ$ for $i=j$.
The script verifies:
\begin{itemize}
\item JEPD, inverse, strong-$\EQ$, RCC5 composition (frame
conditions);
\item the round-7 incidence-tagged label rule keeps every distinct
lap pair $(u_i,u_j)$ at label $\PP$ or $\PPI$, \emph{never} $\EQ$;
\item the round-6 collapsed rule
$\lab(u,v)=L_\Q(q(u),q(v))=L_\Q(\pi,\pi)=\EQ$ does mis-equate
distinct laps (explicit demonstration of the defect).
\end{itemize}
$C_\uparrow$ has no finite model, so SAT is not asserted at the truncation;
the script exercises the architectural claim about the regular
certificate.

\subsection{WP9: split copies for incomparable proper superparts}
\label{ssec:wp9}

\texttt{wp9\_round7\_split\_copies.py} builds the three-element
abstract RCC5 frame $\Delta=\{x,a_\text{sup},b_\text{sup}\}$ with
$\rho(x,a_\text{sup})=\PP$, $\rho(x,b_\text{sup})=\PP$,
$\rho(a_\text{sup},b_\text{sup})=\PO$ (forced by strong-$\EQ$
semantics: $a_\text{sup}$ and $b_\text{sup}$ both contain $x$, hence
share a point, hence $\PO$), $A=\{a_\text{sup}\}$, $B=\{b_\text{sup}\}$,
and verifies:
\begin{itemize}
\item the quotient frame respects JEPD, inverse, strong-$\EQ$, and
RCC5 composition;
\item the pre-quotient occurrence structure
$(x,x_a,x_b,a_\text{sup},b_\text{sup})$ with explicit equality-port
chain $x\eqtag x_a\eqtag x_b$ satisfies the round-7 typed equality
congruence (the three mates have identical $\ell_\Q$-labels to every
outside position);
\item no distinct pair receives the $\self$ tag (round-7
no-self-collapse);
\item $C_{\rm split}:=\exists\PP.A\sqcap\exists\PP.B$ is SAT at $x$.
\end{itemize}
This exercises the round-7 split-copy architecture for multiple
selected proper superparts.

\subsection{Summary}

All three scripts pass.  The empirical layer corroborates the
round-7 architectural claims for the three canonical stress concepts
without relying on any project-internal Python module.

%% ==================================================================
\section{No-automata route vs.\ automata route}
\label{sec:two-routes}
%% ==================================================================

GPT-5.5 has supplied two round-7 proofs of the same decidability
theorem.

\paragraph{No-automata route (canonical).}
\cite{GPT55Round7Sketch} (14 pages) and the expanded companion
\cite{GPT55Round7Expanded} (28 pages) use direct finite-index blocking
with explicit request-closed cycles.  No B\"uchi or parity automata
appear.  This route is the canonical statement because it stays inside
the same conceptual layer as the certificate validity check:
everything is a finite syntactic test on finite alphabets.

\paragraph{Automata route (independent witness).}
\cite{GPT55AutomataRoute} (21 pages) reduces satisfiability to
non-emptiness of a parity tree automaton over the split-forest
alphabet.  This route is an independent reproof of the decidability
theorem using a different finite-state machinery.  It is kept in the
repository as a corroboration of the no-automata route, not as the
canonical statement.

\paragraph{Why this matters.}
The two routes converge on the same theorem (concept satisfiability
of $\ALCIRCC{5}$ over strong abstract atomic RCC5 frames is
decidable).  The no-automata route is a more direct argument
($\omega$-regularity is replaced by a finite combinatorial
cycle-closure condition) and aligns with the verification layer's
preference for finite enumeration.  The automata route is a
classically-styled witness:\@ if a reviewer wants to translate the
proof to a familiar automata-theoretic setting, the parity-automaton
construction is already available.

%% ==================================================================
\section{Status of Claude-authored manuscripts}
\label{sec:status}
%% ==================================================================

The following Claude-authored manuscripts are now superseded by
GPT-5.5's round-7 sketch and expanded companion:

\begin{itemize}
\item \texttt{papers/gpt5.5\_round2/repaired\_split\_forest\_no\_automata\_proof\_v2.tex}
(May 2026, 15 pages, the round-5 (M6$'$) delta);
\item \texttt{papers/gpt5.5\_round2/repaired\_split\_forest\_no\_automata\_proof\_items4to8.tex}
(May 2026, 12 pages, the round-5 items 4--8 delta);
\item \texttt{papers/gpt5.5\_round2/repaired\_split\_forest\_no\_automata\_proof\_v2\_consolidated.tex}
(May 2026, 36 pages, the round-5 consolidation);
\item \texttt{papers/gpt5.5\_round2/repaired\_split\_forest\_no\_automata\_proof\_v2\_consolidated\_round6.tex}
(May 2026, 39 pages, the round-6 internal-coherence pass).
\end{itemize}

All four are retained in the repository as the historical audit
trail.  The conceptual contributions that survive into round-7 are:
\begin{enumerate}
\item the recognition that two $\DR/\PO$ side witnesses of the same
source can be $\PP/\PPI$-comparable (round-5 stress concept, now
worked example in \cite{GPT55Round7Sketch} \S6 and verified in
\S\ref{ssec:wp7}); and
\item the broader two-layer architecture (theoretical proof +
self-contained executable verification artefacts) that surfaced the
round-6 defect.
\end{enumerate}

The round-7 sketch is the new canonical pillar of the decidability
argument.  Claude's role going forward is verification, audit, and
documentation, not authorship of the proof manuscript.

%% ==================================================================
\section{Conclusion}
%% ==================================================================

GPT-5.5's round-7 architectural move (occurrence-sensitive pair
shapes with incidence tags) resolves the central defect of Claude's
round-6 manuscript and absorbs the smaller round-6 (G1)--(G6) and
round-7 (D1)--(D5) issues.  The three self-contained verification
scripts (WP7, WP8, WP9) corroborate the three central stress concepts
of the round-7 architecture.  We adopt GPT-5.5's round-7 sketch as
the new canonical statement; the expanded full-details version is
its detailed companion; the alternative automata-route proof is
retained as an independent witness.  Claude's previous round-5 and
round-6 manuscripts are marked superseded and preserved as the
audit trail.

%% ==================================================================
\begin{thebibliography}{9}

\bibitem{GPT55Round7Review}
GPT-5.5 (OpenAI), prompted by Michael Wessel.
\emph{Formal Review of the Round-6 Consolidated Generated
Split-Forest Proof}.  Unpublished manuscript, May~27, 2026.
Repository:
\url{https://github.com/lambdamikel/alcircc5/blob/master/papers/gpt5.5_final/formal_response_round6_review.pdf}.

\bibitem{GPT55Round7Sketch}
GPT-5.5 (OpenAI), prompted by Michael Wessel.
\emph{A Repaired Split-Forest Decidability Proof for $\ALCIRCC{5}$
Without Automata}.  Unpublished manuscript (consolidated round-7,
14~pages), May~2026.  Repository:
\url{https://github.com/lambdamikel/alcircc5/blob/master/papers/gpt5.5_final/repaired_split_forest_all_in_one_round7.pdf}.

\bibitem{GPT55Round7Expanded}
GPT-5.5 (OpenAI), prompted by Michael Wessel.
\emph{Expanded Generated Split-Forest Decidability Proof for
$\ALCIRCC{5}$: Detailed Companion}.  Unpublished manuscript
(round-7, 28~pages), May~2026.  Repository:
\url{https://github.com/lambdamikel/alcircc5/blob/master/papers/gpt5.5_final/expanded_split_forest_full_details_proof.pdf}.

\bibitem{GPT55AutomataRoute}
GPT-5.5 (OpenAI), prompted by Michael Wessel.
\emph{A Split-Forest Decidability Proof for $\ALCIRCC{5}$ via Parity
Tree Automata: Detailed Version}.  Unpublished manuscript (round-7,
21~pages), May~2026.  Repository:
\url{https://github.com/lambdamikel/alcircc5/blob/master/papers/gpt5.5_final/split_forest_automata_decidability_proof_detailed.pdf}.

\bibitem{ALCIRCC5Repo}
Michael Wessel et al.
\emph{ALCI\_RCC5 Project Repository}.
\url{https://github.com/lambdamikel/alcircc5}.

\end{thebibliography}

\end{document}
